\subsection{Three ways of using MAS}
\label{sec:uses-of-mas}
\slidesrc{Intro to agents \& MAS: 18--41}

Agents and MAS can enter a software project at three different levels:

\begin{center}
\begin{tikzpicture}[font=\small, >=Stealth,
  b/.style={draw, rounded corners, thick, align=center, text width=4.0cm, minimum height=2.0cm}]
  \node[b, draw=nbblue, fill=nbblue!7] (se) at (0,0) {\textbf{1. SE metaphor}\\analysis and design\\think of the system as agents\\[2pt]\footnotesize PDT, INGENIAS};
  \node[b, draw=nbpurple, fill=nbpurple!7] (dev) at (4.7,0) {\textbf{2. Development}\\implementation\\the code \emph{is} agents\\[2pt]\footnotesize JADE, Jason};
  \node[b, draw=nbgreen, fill=nbgreen!7] (sim) at (9.4,0) {\textbf{3. Simulation}\\model a real system\\as interacting agents\\[2pt]\footnotesize NetLogo};
\end{tikzpicture}
\captionof{figure}{The three uses of MAS (slides 18--41). They can be combined: e.g.\ design with
agents and implement with mainstream technologies.}
\end{center}

\subsubsection{MAS as a software engineering metaphor}
\slidesrc{Intro to agents \& MAS: 18--21}

\begin{nbfact}{Zambonelli and Omicini (JAAMAS, 2004)}
``\ldots today's software engineering approaches are increasingly adopting abstractions
approaching that of agent-based computing. This trend can be better understood by
recognising that \emph{the vast majority of modern distributed systems scenarios are
intrinsically prone to be developed in terms of MASs}, and that modern distributed systems
are already de facto MASs, i.e., they are indeed composed of autonomous, situated, and
social components.''
\end{nbfact}

In other words: even if you never write a line of ``agent code'', thinking in terms of
agents (who are the autonomous components? what are their goals? how do they interact?)
helps in the \emph{analysis} and \emph{design} of distributed systems. Tools exist to support
this: \textbf{PDT} (Prometheus Design Tool, RMIT University) and \textbf{INGENIAS}, which
let you draw agents, goals, messages and data stores as diagrams. Methodologies are the
topic of Section~\ref{sec:engineering-mas}.

\subsubsection{Agent-based development}
\slidesrc{Intro to agents \& MAS: 22--36}

\begin{nbdef}{Agent-based development (slides 22--23)}
Agents are employed as the technology for \textbf{implementing the software application
itself}. This choice is usually made when
\begin{itemize}[nosep]
  \item the agent metaphor was already applied in the earlier engineering steps (analysis,
        design);
  \item the intelligence of the \emph{individual} agent is more relevant than the
        ``emerging intelligence'' of the overall system;
  \item complex patterns of interaction may be required;
  \item the application is geographically \textbf{distributed}, \textbf{decentralised}, or
        \textbf{open}.
\end{itemize}
\end{nbdef}

Two agent platforms are shown: \textbf{JADE} (Java Agent DEvelopment framework,
Section~\ref{sec:jade}), whose sniffer tool draws the messages exchanged by agents as a
sequence diagram (slide 24), and \textbf{Jason} (slide 25, an interpreter of AgentSpeak, see
the book \emph{Programming multi-agent systems in AgentSpeak using Jason};
Section~\ref{sec:bdi-agentspeak}). (Slide 25 is labelled ``Jade'' but refers to Jason.)

\begin{nbwarning}[Why industry rarely uses agent platforms (slides 26--27)]
Many agent-based development tools exist, but they are \emph{not commercial projects
supported by large and solid companies}. This is one of the reasons for the limited adoption
of this technology by industry. A pragmatic answer: use technologies that are not
officially tagged as ``agent-oriented'', but are very suitable for implementing software
agents, in the \emph{implementation} stage, and a ``pure'' agent approach for the
\emph{modelling and design} stages.
\end{nbwarning}

The slides then list research works that follow this idea, mapping agents onto mainstream
technologies over the years:

\begin{center}
\small
\begin{tabularx}{\linewidth}{@{}l l X@{}}
\toprule
\textbf{Year} & \textbf{Technology} & \textbf{Work shown on the slide} \\
\midrule
2009 & service-oriented computing & \emph{Agents and Service-Oriented Computing for Autonomic Computing: A Research Agenda} \\
2011 & cloud & \emph{Software Agents as Cloud Computing Services}; \emph{Cloud Computing and Software Agents: Towards Cloud Intelligent Services} \\
2013 & web services & JADE Web Services Integration Gateway (WSIG) \\
2013 & Internet of Things & \emph{From Internet of Things to Internet of Agents} \\
2015 & Internet of Things & \emph{Agents of Things (AoT): an intelligent operational concept of the IoT} \\
2018 & IoT and the Web & EMAS 2018 keynote: \emph{Autonomous Agents for Flexible Hypermedia Systems} \\
2023 & virtual reality & \emph{VEsNA}: virtual environments via natural language agents, applied to factory automation (Gatti, Mascardi) \\
2024 & extended reality & XRIA 2024: first workshop on ``eXtended Reality \& Intelligent Agents'' at ECAI \\
\bottomrule
\end{tabularx}
\captionof{table}{Agents implemented with mainstream technologies (slides 28--36).}
\end{center}

\begin{nbwhy}[Agents everywhere, under other names]
A microservice that watches a queue, decides and calls other services is an agent in
disguise; so is an IoT device that negotiates energy prices, or a non-player character in a
VR training environment that understands spoken commands. The agent \emph{concepts}
(autonomy, goals, interaction) transfer to these technologies even when no agent platform is
used.
\end{nbwhy}

\subsubsection{MAS for simulation}
\slidesrc{Intro to agents \& MAS: 37--41}

\begin{nbfact}{Agent-based modelling and simulation (Macal and North, 2007)}
``Agent-based modeling and simulation (ABMS) is a new approach to modeling systems
comprised of autonomous, interacting agents. ABMS promises to have far-reaching effects on
the way that businesses use computers to support decision-making and researchers use
electronic laboratories to support their research.''

``Some have gone so far as to contend that ABMS is a \emph{third way of doing science}, in
addition to traditional deductive and inductive reasoning. [\ldots] Applications range from
modeling agent behavior in the stock market, supply chains, and consumer markets, to
predicting the spread of epidemics, the threat of bio-warfare, and the factors responsible
for the fall of ancient civilizations.''
\end{nbfact}

The tool shown on slide 41 is \textbf{NetLogo}, a programmable environment where thousands
of simple agents (``turtles'') move on a grid of ``patches'' and the global behaviour is
observed in real time. ABMS is the topic of Section~\ref{sec:abms}.

\begin{nbintuition}[Deduction, induction, simulation]
Deduction: from general laws to particular consequences (logic, Part~II). Induction: from
many observations to general laws (statistics, machine learning). Simulation: write the
\emph{local} rules of each individual, run them, and \emph{observe} what global phenomenon
emerges (a traffic jam, an epidemic wave, a market crash). That is the ``third way''.
\end{nbintuition}
